\documentclass[12pt]{amsart}
\usepackage{amscd,amsmath,amsthm,amssymb}
\usepackage{mathtools}
\usepackage[margin=2cm]{geometry}
\usepackage{epsfig}
\usepackage{rawfonts}
\usepackage{enumerate}
\usepackage{graphics}
\usepackage{multirow}
\usepackage{xspace}
\usepackage{graphicx}
\usepackage{mathrsfs}
\usepackage{amsmath}
\usepackage{amsfonts}
\usepackage{amssymb}
\usepackage{amsthm}
\usepackage{graphicx}
\usepackage{booktabs}
\usepackage{caption}
\usepackage{listings}
\usepackage{setspace}
\usepackage[mathscr]{eucal}
\usepackage{pgfplots}
\usepackage{hyperref}
\usepackage{wrapfig}
\usepackage{floatflt,epsfig}
\usepackage{ dsfont }
\usepackage{amscd}
\usepackage{tikz-cd}
\usepackage{fancyhdr}
\usepackage[all]{xy}
\usepackage{latexsym}
\usepackage{amscd}
\usepackage{pifont}
\usepackage{subfig}
\usepackage{easyReview}
\usepackage{subfig}
\usepackage{pstricks-add}
\usepackage{pgf,tikz,pgfplots}
\pgfplotsset{compat=1.15}
\usepackage{mathrsfs}
\usetikzlibrary{arrows}
\usetikzlibrary[patterns]
%\usepackage{subcaption}
 \usepackage[normalem]{ulem}
 \usepackage{multicol}

 
%
%------    GENERAL MACROS    -----
%
% Standard rings and fields, affine and projective space
%
\def\NZQ{\Bbb}               % the font for N,Z,Q,R,C
\def\NN{{\NZQ N}}
\def\QQ{{\NZQ Q}}
\def\ZZ{{\NZQ Z}}
\def\RR{{\NZQ R}}
\def\CC{{\NZQ C}}
\def\AA{{\NZQ A}}
\def\PP{{\NZQ P}}
\def\DD{{\NZQ D}}
\def\FF{{\NZQ F}}
\def\GG{{\NZQ G}}
\def\HH{{\NZQ H}}
\def\EE{{\NZQ E}}
\def\KK{{\NZQ K}}
\def\MM{{\NZQ M}}
%
%------------------------------------------------
% Symbols in "Fraktur"
%
\def\frk{\frak}               % font for "Fraktur"
\def\aa{{\frk a}}
\def\pp{{\frk p}}
\def\Pp{{\frk P}}
\def\qq{{\frk q}}
\def\Qq{{\frk Q}}
\def\mm{{\frk m}}
\def\Mm{{\frk M}}
\def\nn{{\frk n}}
%\def\Phi{{\frk N}}
%
%------------------------------------------------
\def\Ic{{\mathcal I}}
\def\Jc{{\mathcal J}}
\def\Rc{{\mathcal R}}
\def\Lc{{\mathcal L}}
\def\Gc{{\mathcal G}}
\def\Fc{{\mathcal F}}
\def\Dc{{\mathcal D}}
\def\Pc{{\mathcal P}}
\def\Qc{{\mathcal Q}}
\def\Mc{{\mathcal M}}
\def\Hc{{\mathcal H}}
\def\Gc{{\mathcal G}}
\def\Ac{{\mathcal A}}
\def\Bc{{\mathcal B}}
\def\Cc{{\mathcal C}}
\def\Wc{{\mathcal W}}
\def\Sc{{\mathcal S}}
\def\Fc{{\mathcal F}}

\renewcommand{\qedsymbol}{$\square$}
\newcommand{\rev}{\mathrm{rev}}


% Small letters in bold
%
\def\ab{{\bold a}}
\def\bb{{\bold b}}
\def\xb{{\bold x}}
\def\yb{{\bold y}}
\def\zb{{\bold z}}
\def\cb{{\bold c}}
\def\db{{\bold d}}
\def\eb{{\bold e}}
%
\def\opn#1#2{\def#1{\operatorname{#2}}} % to make operators
%------------------------------------------------
% Numerical invariants of rings, ideals, and modules
%
\opn\chara{char} \opn\length{\ell} \opn\pd{pd} \opn\rk{rk}
\opn\projdim{proj\,dim} \opn\injdim{inj\,dim} \opn\rank{rank}
\opn\depth{depth} \opn\grade{grade} \opn\height{height}
\opn\embdim{emb\,dim} \opn\codim{codim}
\def\OO{{\mathcal O}}
\opn\Tr{Tr} \opn\bigrank{big\,rank}
\opn\superheight{superheight}\opn\lcm{lcm}
\opn\trdeg{tr\,deg}%\emph{
	\opn\reg{reg} \opn\lreg{lreg} \opn\ini{in} \opn\lpd{lpd}
	\opn\size{size} \opn\sdepth{sdepth}
	\opn\link{link}\opn\fdepth{fdepth}\opn\lex{lex}\opn\dist{dist}
	%------------------------------------------------
	% Divisors
	%
	\opn\div{div} \opn\Div{Div} \opn\cl{cl} \opn\Cl{Cl}
	%
	%------------------------------------------------
	% Subsets of the spectrum of a ring
	%
	\opn\Spec{Spec} \opn\Supp{Supp} \opn\supp{supp} \opn\Sing{Sing}
	\opn\Ass{Ass} \opn\Min{Min}\opn\Mon{Mon}
	%
	%------------------------------------------------
	% Standard operations on ideals and modules
	%
	\opn\Ann{Ann} \opn\Rad{Rad} \opn\Soc{Soc}
	%
	%------------------------------------------------
	% Linear algebra and homology, endo- and automorphisms
	%
	\opn\Im{Im} \opn\Ker{Ker} \opn\Coker{Coker} \opn\Am{Am}
	\opn\Hom{Hom} \opn\Tor{Tor} \opn\Ext{Ext} \opn\End{End}
	\opn\Aut{Aut} \opn\id{id}
	\def\Frob{{\mathcal F}}
	\opn\nat{nat}
	\opn\pff{pf}%   \pf exists already
	\opn\Pf{Pf} \opn\GL{GL} \opn\SL{SL} \opn\mod{mod} \opn\ord{ord}
	\opn\Gin{Gin} \opn\Hilb{Hilb}\opn\sort{sort}
	%
	%------------------------------------------------
	% Convexity
	%
	\opn\aff{aff} \opn
	\con{conv} \opn\relint{relint} \opn\st{st}
	\opn\lk{lk} \opn\cn{cn} \opn\core{core} \opn\vol{vol}
	\opn\link{link} \opn\star{star}\opn\lex{lex}\opn\set{set}
	%------------------------------------------------
	% Graded rings and Rees algebras
	\opn\gr{gr}
	\def\Rees{{\mathcal R}}
	%
	%------------------------------------------------
	% Polynomials and power series
	%
	\def\poly#1#2#3{#1[#2_1,\dots,#2_{#3}]}
	\def\pot#1#2{#1[\kern-0.28ex[#2]\kern-0.28ex]}
	\def\Pot#1#2#3{\pot{#1}{#2_1,\dots,#2_{#3}}}
	\def\konv#1#2{#1\langle #2\rangle}
	\def\Konv#1#2#3{\konv{#1}{#2_1,\dots,#2_{#3}}}
	%
	%------------------------------------------------
	% Direct and inverse limits
	%
	\opn\dirlim{\underrightarrow{\lim}}
	\opn\inivlim{\underleftarrow{\lim}}
	%
	%
	% Names with a meaning
	%
	\let\union=\cup
	\let\sect=\cap
	\let\dirsum=\oplus
	\let\tensor=\otimes
	\let\iso=\cong
	\let\Union=\bigcup
	\let\Sect=\bigcap
	\let\Dirsum=\bigoplus
	\let\Tensor=\bigotimes
	%
	%------------------------------------------------
	%
	\let\to=\rightarrow
	\let\To=\longrightarrow
	\def\Implies{\ifmmode\Longrightarrow \else
		\unskip${}\Longrightarrow{}$\ignorespaces\fi}
	\def\implies{\ifmmode\Rightarrow \else
		\unskip${}\Rightarrow{}$\ignorespaces\fi}
	\def\iff{\ifmmode\Longleftrightarrow \else
		\unskip${}\Longleftrightarrow{}$\ignorespaces\fi}
	\let\gets=\leftarrow
	\let\Gets=\longleftarrow
	\let\followsfrom=\Leftarrow
	\let\Followsfrom=\Longleftarrow
	\let\:=\colon
	%
	%
	%
	%We like the var forms of some greek letters (as taught in German schools)
	\let\epsilon\varepsilon
	%\let\phi=\varphi
	\let\kappa=\varkappa
 %
	% %           We print on A4 paper
	% %
	% \textwidth=15cm \textheight=22cm \topmargin=0.5cm
	% \oddsidemargin=0.5cm \evensidemargin=0.5cm \pagestyle{plain}
        % %
	% %           The pf environment of AMSART needs a little help
	% %
	\def\qed{\ifhmode\textqed\fi
		\ifmmode\ifinner\quad\qedsymbol\else\dispqed\fi\fi}
	\def\textqed{\unskip\nobreak\penalty50
		\hskip2em\hbox{}\nobreak\hfil\qedsymbol
		\parfillskip=0pt \finalhyphendemerits=0}
	\def\dispqed{\rlap{\qquad\qedsymbol}}
	\def\noqed{\def\qed{\relax}}
	%
	% ------    END OF GENERAL MACROS    -------
	\opn\dis{dis}
	\def\pnt{{\raise0.5mm\hbox{\large\bf.}}}
	\def\lpnt{{\hbox{\large\bf.}}}
	\opn\Lex{Lex}
	
	%-- macro for local cohomology-----------------------------
	\def\Coh#1#2{H_{\mm}^{#1}(#2)}
	%-- macro for a complicated condition for the extended
	%-- Hochster's formula
	\def\hchst#1{for all $u\in G(I)$ there  exists $i\notin {#1}$
		such that $\nu_i(u) > a_i \geq 0$}	



    %%%%%%%%%%%%% Enumeration Theorems, Propositions, and so on %%%%%%%%%%%%%%%%%%%%%%%%%%
        
        \newtheorem{Theorem}{Theorem}[section]
	\newtheorem{Lemma}[Theorem]{Lemma}
	\newtheorem{Corollary}[Theorem]{Corollary}
	\newtheorem{Proposition}[Theorem]{Proposition}
	\newtheorem{Remark}[Theorem]{Remark}
	\newtheorem{Remarks}[Theorem]{Remarks}
	\newtheorem{Example}[Theorem]{Example}
	\newtheorem{Examples}[Theorem]{Examples}
	\newtheorem{Definition}[Theorem]{Definition}
	\newtheorem{Problem}[Theorem]{}
	\newtheorem{Conjecture}[Theorem]{Conjecture}
	\newtheorem{Question}[Theorem]{Question}
        \newtheorem{Discussion}[Theorem]{Discussion}
        \newtheorem{Notation}[Theorem]{Notation}
                \newtheorem{Assumption}[Theorem]{Assumption}
       \newtheorem{Setup}[Theorem]{Setup}
       \newtheorem{Construction}[Theorem]{Construction}

\begin{document}
   
        \title[ideals of $2$-minors]{Prime and Cohen--Macaulay ideals of $2$-minors with linear resolution}
	
\author[T.~Hibi]{Takayuki Hibi}
\address[Takayuki Hibi]
{Department of Pure and Applied Mathematics, 
Graduate School of Information Science and Technology, 
Osaka University, 
Suita, Osaka 565-0871, Japan}
\email{hibi@math.sci.osaka-u.ac.jp}

\author[A. A. Qureshi]{Ayesha Asloob Qureshi}      
   \address[Ayesha Asloob Qureshi]{Sabanci University, Faculty of Engineering and Natural Sciences, Orta Mahalle, Tuzla 34956, Istanbul, Turkey}	
%mysu
\email{aqureshi@sabanciuniv.edu, ayesha.asloob@sabanciuniv.edu}

% \author[S.~Saeedi~Madani]{Sara Saeedi Madani}
% \address[Sara Saeedi Madani]
% {Department of Mathematics and Computer Science, Amirkabir University of Technology, Tehran, Iran, and School of Mathematics, Institute for Research in Fundamental Sciences, Tehran, Iran} 
% \email{sarasaeedi@aut.ac.ir, sarasaeedim@gmail.com}
       
	\keywords{prime ideal, Cohen--Macaulay ring, linear resolution}
	
	% \subjclass[2010]{13D02, 05E40, 05E45, 05B50, 05C70}
    \subjclass[2020]{13H10, 13P10}
    	  %05C31, graph polynomials, 	13P25  	
    
	\thanks{} 

        \maketitle
      
\begin{abstract}
Prime and Cohen--Macaulay ideals of $2$-minors of matrices of variables with linear resolution are classified.
\end{abstract}





%=%=%=%=%=%=%=%=%=%=%=%=%=%=%=%=%=%=%=%=%=%=%=%=%=%=%=%=%=%=%=%=%=%=%=%=%=%=%=%=%=%=%=%=%=%=%=%=%=%=%=%=%=%=%=%=%=%=%=%=%=%=%=%=%=%=%=%=%=%=%=%=


\section*{Introduction}
Binomial ideals form an important class of ideals in commutative algebra,
with close connections to combinatorics and toric geometry. Among the
classical examples are determinantal ideals generated by the $2$-minors
of a matrix of variables. In particular, let
$S=K[x_1,\ldots,x_n,y_1,\ldots,y_n]$ be a polynomial ring over a field
$K$, and let
\[
X_0=
\begin{bmatrix}
x_1 & x_2 & \cdots & x_n\\
y_1 & y_2 & \cdots & y_n
\end{bmatrix}
\]
be the generic $2\times n$ matrix. Set
\[
J=I_2(X_0)
  =(x_i y_j-x_jy_i:1\leq i<j\leq n).
\]
It is well known that $J$ is prime, $S/J$ is Cohen--Macaulay, and $J$
has a linear resolution; see, for example, \cite{BrunsVetter}.

In our previous work \cite{HQS}, we studied quadratic binomial ideals
from a converse point of view. Namely, we asked to what extent a prime
quadratic binomial ideal with Cohen--Macaulay quotient and linear
resolution must arise as the ideal of $2$-minors of a matrix of
variables. This question was motivated in part by several families of
binomial ideals arising from combinatorial structures for which these
three properties lead naturally to determinantal ideals. In
\cite{HQS}, we obtained partial results in this direction and were led
to the following related classification problem.

Let $X$ be a $2\times n$ matrix of not necessarily distinct variables in $S$. Let $I_2(X)\subseteq S$ denote the ideal generated by the $2$-minors of
$X$. We call $X$ a {\em PCML} matrix if (i) $I_2(X)$ is prime, (ii) $S/I_2(X)$ is Cohen--Macaulay and (iii) $I_2(X)$ has linear resolution. The problem of classifying all PCML matrices was posed in \cite{HQS}.
The purpose of the present paper is to solve this problem for
$2\times n$ matrices.

After excluding the immediate obstructions to primeness, we may write
\[
X=
\begin{bmatrix}
x_1 & x_2 & \cdots & x_n\\
z_1 & z_2 & \cdots & z_n
\end{bmatrix},
\]
where $z_1,\ldots,z_n$ are pairwise distinct elements of
$\{x_1,\ldots,x_n,y_1,\ldots,y_n\}$ and $z_i\neq x_i$ for every $i$.
Allowing repetitions between the two rows makes the classification
considerably subtler than in the generic case. Nevertheless, there are
important nongeneric examples retaining the PCML property, including
the $2\times n$ Hankel matrix. Thus it is natural to ask which patterns
of repetitions preserve these algebraic and homological properties.

In Section~1, for each subset
$\delta\subseteq\{2,\ldots,n\}$, we introduce the matrix
\[
X_n(\delta)=
\begin{bmatrix}
x_1 & x_2 & \cdots & x_n\\
\xi_1 & \xi_2 & \cdots & \xi_n
\end{bmatrix},
\qquad
\xi_i=
\begin{cases}
y_i, & i\notin\delta,\\
x_{i-1}, & i\in\delta.
\end{cases}
\]
This family contains both the generic matrix, when $\delta=\emptyset$,
and the Hankel matrix, when $\delta=\{2,\ldots,n\}$. In
Theorem~\ref{prime+CM+linear}, we prove that every $X_n(\delta)$ is PCML.

In Section~2, we associate with $X$ a directed graph $G_X$ on the
vertex set $\{1,\ldots,n\}$ by declaring that
\[
i\longrightarrow j
\qquad\Longleftrightarrow\qquad
x_i=z_j.
\]
As observed in Remark~\ref{rem:def}, every connected component of $G_X$ is either
a directed path or a directed cycle. In Theorem~\ref{prime}, we prove that
$I_2(X)$ is prime if and only if $G_X$ contains no directed cycle.
Moreover, in the acyclic case, the directed paths can be used to
permute the columns and relabel the variables so that $X$ becomes a
matrix of the form $X_n(\delta)$. Combining this with Theorem~\ref{prime+CM+linear}
yields our main classification, Theorem~\ref{mainPCML}:
\[
X\text{ is PCML}
\quad\Longleftrightarrow\quad
I_2(X)\text{ is prime}
\quad\Longleftrightarrow\quad
G_X\text{ contains no directed cycle}.
\]
Thus, within this class of $2\times n$ matrices, primeness alone forces
the Cohen--Macaulay property and linear resolution.

Theorem~\ref{mainPCML} also naturally raises the question whether the remaining two properties can be characterized combinatorially without assuming primality. More precisely, for a matrix $X$ as above, it would be interesting to determine, in terms of the directed graph $G_X$, when $S/I_2(X)$ is Cohen--Macaulay and when $I_2(X)$ has a linear resolution. In particular, one may ask how the structure of the directed cycles of $G_X$ influences these two properties. A broader direction is to ask whether an analogous classification can be obtained for $m\times n$ matrices with $m>2$. In that setting, the possible patterns of repeated variables are substantially more complicated, and the simple directed-graph description used here is no longer immediately available. Understanding to what extent the classification obtained in the present paper admits a higher-dimensional analogue therefore appears to be a natural problem for further study.




%=%=%=%=%=%=%=%=%=%=%=%=%=%=%=%=%=%=%=%=%=%=%=%=%=%=%=%=%=%=%=%=%=%=%=%=%=%=%=%=%=%=%=%=%=%=%=%=%=%=%=%=%=%=


\section{Ideals of $2$-minors of $2 \times n$ matrices}
Let $S=K[x_1,\ldots,x_n,y_1,\ldots,y_n]$ denote the polynomial ring in $2n$ variables over a field $K$.  Fix $\delta \subset \{2, \ldots, n\}$ and introduce a $2 \times n$ matrix
$$
X_n(\delta)=
\begin{bmatrix}
   x_1 & x_2 & \cdots & x_n \\
   \xi_{1} & \xi_{2} & \cdots & \xi_{n} 
\end{bmatrix}
$$
where 
\[
\xi_i
=
\begin{cases}
\,y_i, & \text{if} \quad i \not\in \delta,\\[2pt]
\,x_{i-1},     & \text{if} \quad i \in \delta.
\end{cases}
\]
In particular, 
\[
X_n(\{2,\ldots, n\}) \, = \, 
\begin{bmatrix}
   x_1 & x_2 & x_3 & \cdots & x_n \\
   y_{1} & x_1 & x_2 & \cdots & x_{n-1}  
\end{bmatrix},
\qquad
X_n(\emptyset)=
\begin{bmatrix}
   x_1 & x_2 & \cdots & x_n \\
   y_{1} & y_{2} & \cdots & y_n  
\end{bmatrix}.
\]
Let $S_\delta$ denote the polynomial ring in the variables $x_1, \ldots, x_n$ and $y_j$ with $j \not\in \delta$ over $K$ and $I_n(\delta)$ the ideal of $S_\delta$ generated by all $2$-minors of $X_n(\delta)$. 


The two extreme cases of this family are classical. If
$\delta=\emptyset$, then $I_n(\delta)$ is the ideal of maximal minors
of the generic $2\times n$ matrix. It is well known that this ideal is
prime, its quotient ring is Cohen--Macaulay, and it has a linear
resolution; see, for example, \cite{BrunsVetter}. At the other extreme,
if $\delta=\{2,\ldots,n\}$, then, after relabeling the variables,
$I_n(\delta)$ is the ideal generated by the $2$-minors of a
$2\times n$ Hankel matrix. This is the defining ideal of a rational
normal curve; in particular, it is prime, its quotient ring is
Cohen--Macaulay, and its minimal free resolution is given by the
Eagon--Northcott complex, hence the ideal has a linear resolution;
see \cite{ConcaRNC}.

The following theorem shows that these three properties persist for
every intermediate choice of
$\delta\subseteq\{2,\ldots,n\}$, thereby interpolating between the
generic and Hankel cases.

\begin{Theorem}\label{prime+CM+linear}
One has:
\begin{enumerate}
\item[(i)] the ideal $I_n(\delta)\subset S_\delta$ is prime;
\item[(ii)] the quotient ring $S_\delta/I_n(\delta)$ is Cohen--Macaulay;
\item[(iii)] $I_n(\delta)$ has a linear resolution.
\end{enumerate}
\end{Theorem}

We begin by proving parts (ii) and (iii) of
Theorem~\ref{prime+CM+linear}. The idea is to obtain
$I_n(\delta)$ from the generic determinantal ideal $I_n(\emptyset)$
by imposing the linear relations
\[
x_{i-1}-y_i,\qquad i\in\delta.
\]
We first record the elementary identities that make this reduction precise. Let $i\notin\delta$ and $j\in\delta$. Then
\[
(x_i x_{j-1}-x_jy_i,\,x_{j-1}-y_j)
=
(x_i y_j-x_jy_i,\,x_{j-1}-y_j).
\]
Similarly, if $i,j\notin\delta$ with $i\neq j$, then
\begin{eqnarray*}
&&
(x_i x_{j-1}-x_jx_{i-1},\,
  x_{i-1}-y_i,\,
  x_{j-1}-y_j)
\\
&=&
(x_i y_j-x_jy_i,\,
  x_{i-1}-y_i,\,
  x_{j-1}-y_j).
\end{eqnarray*}
It follows that
\[
S/(I_n(\emptyset),\{x_{i-1}-y_i:i\in\delta\})
=
S/(I_n(\delta),\{x_{i-1}-y_i:i\in\delta\}).
\]

\begin{Lemma}
\label{regular_sequence}
One has
\[
S/(I_n(\emptyset),\{x_{i-1}-y_i:i\in\delta\})
\cong S_\delta/I_n(\delta).
\]
\end{Lemma}

\begin{proof}
Modulo the relations $x_{i-1}-y_i$, $i\in\delta$, each variable
$y_i$ with $i\in\delta$ is identified with $x_{i-1}$. Hence
\[
S/(\{x_{i-1}-y_i:i\in\delta\})\cong S_\delta.
\]
Under this identification, the image of $I_n(\emptyset)$ is precisely
$I_n(\delta)$, by the identities above. The assertion follows.
\end{proof}

We can now deduce the Cohen--Macaulay property and linearity of the
resolution from the corresponding properties of the generic
determinantal ideal.

\begin{Corollary}
\label{CM_linear}
The quotient ring $S_\delta/I_n(\delta)$ is Cohen--Macaulay and
$I_n(\delta)$ has a linear resolution.
\end{Corollary}

\begin{proof}
Recall that $S/I_n(\emptyset)$ is Cohen--Macaulay and
$I_n(\emptyset)$ has a linear resolution. Since
\[
\{x_{i-1}-y_i:i\in\delta\}
\]
is a regular sequence of $S/I_n(\emptyset)$
(\cite[Theorem~4.1]{Eis}), Lemma~\ref{regular_sequence} implies that
$S_\delta/I_n(\delta)$ is Cohen--Macaulay. Furthermore, the $h$-vector of $S_\delta/I_n(\delta)$ coincides with
that of $S/I_n(\emptyset)$. Since $I_n(\emptyset)$ has a linear
resolution, it follows that $I_n(\delta)$ has a linear resolution.
\end{proof}


We now turn to a proof of (i) of Theorem \ref{prime+CM+linear}.
Let $<_{\rm rev}$ denote the reverse lexicographic order of $S$ induced by the ordering of the variables 
\[
y_1 < x_1 < y_2 < x_2 < y_3 < \cdots < y_{n} < x_n
\]
and $<_{\rm rev}^*$ the restriction of $<_{\rm rev}$ on $S_\delta$.  Let ${\rm in}_{<_{\rm rev}^*}(I_n(\delta))$ denote the initial ideal of $I_n(\delta)$ with respect to $<_{\rm rev}^*$.  Let $\Omega=A \cup B$, where 
    \begin{eqnarray*}
&A=\{x_k y_\ell:1 \leq k \leq n, \ell \not\in \delta, k < \ell\},& \\
&B=\{x_k x_{\ell-1}:1 \leq k \leq n, \ell \in \delta, k < \ell\}.&
    \end{eqnarray*}
and $(\Omega) \subset S_\delta$ the monomial ideal generated by those monomials belonging to $\Omega$.
\begin{Lemma}
    One has $(\Omega) \subset {\rm in}_{<_{\rm rev}^*}(I_n(\delta))$.
\end{Lemma}

Suppose that $2 \not\in \delta$ and define the ring homomorphism $$\varphi:S_\delta \to K[y_1,x_1^{-1}, \{x_i: i+1 \not\in \delta\}]$$ by setting
\begin{itemize}
    \item 
    $\varphi(y_1)=y_1$,
    \item 
    $\varphi(x_i)=x_i, \, i+1 \not\in \delta$,
    \item 
    $\varphi(y_j)=(y_1/x_1)^{i_0-j+1}x_{i_0}, \, \text{where $1 < j \not\in \delta$ and $i_0$ is the smallest $i$ with $i+1 \not\in \delta$ and  $j \leq i$}$,
    \item
    $\varphi(x_j)=(y_1/x_1)^{i_0-j}x_{i_0}, \, \text{where $j+1 \in \delta$ and $i_0$ is the smallest $i$ with $i+1 \not\in \delta$ and  $j < i$}$.
\end{itemize}


For example, if $n=7$ and $\delta=\{3,4,6\}$, then
\begin{eqnarray*}
&&\varphi\left(
\begin{bmatrix}
   x_1 & x_2 & x_3 & x_4 & x_5 & x_6 & x_7 \\
   y_1 & y_2 & x_2 & x_3 & y_5 & x_5 & y_7  
\end{bmatrix}
\right) \\
&&=
\begin{bmatrix}
   x_1 & (y_1/x_1)^2x_4 & (y_1/x_1)x_4 & x_4 & (y_1/x_1)x_6 & x_6 & x_7 \\   y_1 & (y_1/x_1)^3x_4 & (y_1/x_1)^2x_4 & (y_1/x_1)x_4 & (y_1/x_1)^2x_6 & (y_1/x_1)x_6 & (y_1/x_1)x_7  
\end{bmatrix}.
\end{eqnarray*}




One has 
\[
I_n(\delta) \subset {\rm Ker}(\varphi), \quad {\rm in}_{<_{\rm rev}^*}(I_n(\delta)) \subset {\rm in}_{<_{\rm rev}^*}({\rm Ker}(\varphi)).
\]
Let $\Mc_\delta$ denote the set of monomials of $S_\delta$ and 
\[
\Omega^\# := \{u \in \Mc_\delta : u \not\in (\Omega)\}.
\]
\begin{Lemma}
    \label{linear_independent}
The set of Laurent monomials $\{\varphi(u) : u \in \Omega^\# \}$ is linearly independent. 
\end{Lemma}

\begin{proof}
Let $\supp{(u)}$ denote the support of a monomial $u \in S_\delta$.  Our task is to prove that if $u,v \in \Omega^\#$ and if $u$ and $v$ are relatively prime, then $\varphi(u) \neq \varphi(v)$.  Set 
\[
V_\delta := \{x_i : i+1 \not\in \delta\} \cup \{y_i : i \not\in \delta\}.
\]
Let
\[
\{x_i : i+1 \not\in \delta\} = \{x_1, x_{i_1}, \ldots, x_{i_{n-|\delta|-1}}\}, \quad 1 < i_1 < \cdots < i_{n-|\delta|-1}
\]
and
\[
\{y_i : i \not\in \delta\} = \{y_1, y_{j_1}, \ldots, y_{j_{n-|\delta|-1}}\}, \quad 1 < j_1 < \cdots < j_{n-|\delta|-1}.
\]
One has 
\[
\varphi(y_{j_k}) = (y_1/x_1)^{i_k - j_k + 1} x_{i_k}, \quad k = 2,\ldots, n-|\delta|-1.
\]

\medskip

\noindent
{\bf (First Step)}  If $\supp(u) \cup \supp(v) \subset V_\delta$, then the desired result follows from \cite[p.~98]{Hibi}.   

% Let $\supp(u) \cap \{x_i : i+1 \not\in \delta\} \neq \emptyset$ and $\supp(v) \cap \{x_i : i+1 \not\in \delta\} \neq \emptyset$.  Let $i_u$ denote the smallest $i$ with $i+1 \not\in \delta$ for which $x_i \in \supp(u)$.  Let $i_u < i_v$.  Let $\varphi(y_{j}) = (y_1/x_1)^{i_v - j + 1} x_{i_v}$.  Since $y_{j} \not\in \supp(u)$, it follows that  $\varphi(u) \neq \varphi(v)$. 

% Let $\supp(u) \cap \{x_i : i+1 \not\in \delta\} \neq \emptyset$ and $\supp(v) \subset \{y_i : i \not\in \delta\}$.  If $(\supp(u) )\cap \{y_i : i \not\in \delta \} \neq \emptyset,$ then clearly $\varphi(u) \neq \varphi(v)$.  If $\supp(u)\cap \{y_i : i \not\in \delta, i < i_u\} = \emptyset,$ then since no $y_1/x_1$ can appear in $\varphi(u)$, one has $\varphi(u) \neq \varphi(v)$. 

\medskip

\noindent
{\bf (Second Step)}  Let $\supp(u) \subset V_\delta$ and $\supp(v) \not\subset V_\delta$.  Let $i_0$ denote the unique integer $i$ for which $x_{i} \in \supp(v)$ and $i+1 \in \delta$.  One has 
\[
\supp(v) \subset \{x_{i_0}\} \cup \{x_i \in V_\delta: i > i_0\} \cup \{y_j \in V_\delta: j \leq i_0 \} .
\]
Let $i^* > i_0$ denote the smallest integer $i$ for which $x_i \in V_\delta$ and $i_* < i_0$ the biggest integer $i$ for which $y_i \in V_\delta$.  
%
Let $x_{i^*} \in \supp(u)$.  Since $y_j \not\in \supp(u)$ for $j > i_0$, it follows that $\varphi(u) \neq \varphi(v)$. 

Let $x_{i^*} \not\in \supp(u)$.  Then $y_{i_*} \in \supp(u)$.  If $\supp(u) \not\subset \{y_i : i \not\in \delta \}$, then clearly $\varphi(u) \neq \varphi(v)$.  If $\supp(u) \subset \{y_i : i \not\in \delta \}$, then considering the powers of $y_1/x_1$, one has $\varphi(u) \neq \varphi(v)$. 

\medskip

\noindent
{\bf (Third Step)}  Let $\supp(u) \not\subset V_\delta$ and $\supp(v) \not\subset V_\delta$.  Let $i_u$ (resp. $i_v$) denote the unique integer $i$ for which $x_{i} \in \supp(u)$ (resp. $x_{i} \in \supp(v)$) and $i+1 \in \delta$.  Let $i_u < i_v$.

If there is no $y_j$ with $i_u < j < i_v$, then, since 
\[
(\supp(u) \cup \supp(v) )\cap V_\delta \subset \{x_i \in V_\delta : i > i_v\} \cup \{y_j : j \leq i_u\}, 
\]
one has $\varphi(u) \neq \varphi(v)$.

Suppose that there is $y_j$ with $i_u < j < i_v$.  Let $\varphi(x_{i_u}) = (y_1/x_1)^a x_{i_*}$ and $\varphi(x_{i_v}) = (y_1/x_1)^b x_{i^*}$, where $i_* < i^*$.  Let $\varphi(y_{j_*}) = (y_1/x_1)^{a'} x_{i_*}$ and $\varphi(y_{j^*}) = (y_1/x_1)^{b'} x_{i^*}$, where $j_* < j^*$.  Now, since $x_{i^*} \in \supp(u)$ and $y_{j_*} \in \supp(v)$,  considering the powers of $y_1/x_1$, one has $\varphi(u) \neq \varphi(v)$. 
\, \, \, 
\end{proof}

When $2 \in \delta$, since $x_1 \not\in V_\delta$, slightly modifying the map $\varphi$ is required.  Let $q \geq 2$ denote the smallest integer $i$ for which $x_i \in V_\delta$.  We then define $\varphi: S_\delta \to K[y_1, x_{q}^{-1}, \{x_i : i+1 \not\in \delta\}]$ by starting with $\varphi(y_1) = y_1^{q}, \varphi(x_1) = x_{q}y_1^{q-1}$ and $\varphi(x_i) = x_i^{q}$ if $x_i \in V_\delta$.  More precisely,  

\begin{itemize}
    \item 
    $\varphi(y_1)=y_1^q$,
    \item 
    $\varphi(x_1) = x_{q}y_1^{q-1}$,
    \item 
    $\varphi(x_i)=x_i^q, \, i+1 \not\in \delta$,
    \item 
    $\varphi(y_j)=(y_1/x_q)^{i_0-j+1}x_{i_0}^q, \, \text{where $1 < j \not\in \delta$ and $i_0$ is the smallest $i$ with $i+1 \not\in \delta$ and  $j \leq i$}$,
    \item
    $\varphi(x_j)=(y_1/x_q)^{i_0-j}x_{i_0}^q, \, \text{where $j+1 \in \delta$ and $i_0$ is the smallest $i$ with $i+1 \not\in \delta$ and  $j < i$}$.
\end{itemize}
Our proof of Lemma \ref{linear_independent} can be valid with slight modifications. 


For example, if $n=5$ and $\delta=\{2,4\}$, then $q=2$ and
\begin{eqnarray*}
&&\varphi\left(
\begin{bmatrix}
   x_1 & x_2 & x_3 & x_4 & x_5  \\
   y_1 & x_1 & y_3 & x_3 & y_5   
\end{bmatrix}
\right) \\
&&=
\begin{bmatrix}
   x_2 y_1& x_2^2 & (y_1/x_2)x_4^2 & x_4^2 & x_5^2  \\   
   y_1^2 & x_2 y_1 & (y_1/x_2)^2x_4^2 & (y_1/x_2)x_4^2 & (y_1/x_2)x_5^2   
\end{bmatrix}.
\end{eqnarray*}



Now, a standard fact on Gr\"obner basis \cite[Proposition 2.2.5]{HHgtm260} guarantees that 

\begin{Corollary}
\label{toric_ideal}
One has $$(\Omega) = {\rm in}_{<_{\rm rev}^*}(I_n(\delta)) = {\rm in}_{<_{\rm rev}^*}({\rm Ker}(\varphi)$$ and $I_n(\delta) = {\rm Ker}(\varphi)$.  In particular, $I_n(\delta)$ is prime.
\end{Corollary}

%%%%%%%%%%%%%%%%%%%%%%%%%%%%%%%%%%%%%%%%%%%%%%%%%%%%%%%%%%%%%%%

\section{Classification of PCML matrices}

Let $X$ be a $2\times n$ matrix of not necessarily distinct variables
belonging to a polynomial ring $S$ over a field $K$, and let $I$ be
the ideal generated by all $2$-minors of $X$. If a variable appears more than once in the same row, or if the two
entries of some column are equal, then $I$ is not prime. Therefore,
in classifying PCML matrices, we may restrict our attention to matrices
whose entries in each row are pairwise distinct and whose two entries
in every column are distinct. After relabeling the variables, we may thus write
\[
X=
\begin{bmatrix}
x_1 & x_2 & \cdots & x_n\\
z_1 & z_2 & \cdots & z_n
\end{bmatrix},
\]
where $z_1,\ldots,z_n$ are pairwise distinct elements of
$\{x_1,\ldots,x_n,y_1,\ldots,y_n\}$ and $z_i\neq x_i$ for
$i=1,\ldots,n$.

Toward the classification of PCML matrices, we first characterize those matrices whose ideals of $2$-minors are prime. For this purpose, we associate a directed graph to $X$. We first recall the graph-theoretic terminology that will be used.

Let $G$ be a directed graph. A \emph{directed path} in $G$ is a sequence of distinct vertices \[ v_1,v_2,\ldots,v_r \] such that $(v_i,v_{i+1})\in E(G)$ for each $i=1,\ldots,r-1$. In particular, an isolated vertex is regarded as a directed path. A \emph{directed cycle} of length $r$ in $G$ consists of a directed path
$v_1,v_2,\ldots,v_r$, with $r\geq 2$, which also satisfies
$(v_r,v_1)\in E(G)$. In particular, if $(u,v)$ and
$(v,u)$ belong to $E(G)$, then $u,v$ form a directed cycle of length $2$.

\begin{Definition}
{\em
Let
\[
X=
\begin{bmatrix}
x_1 & x_2 & \cdots & x_n\\
z_1 & z_2 & \cdots & z_n
\end{bmatrix},
\]
where $z_1,\ldots,z_n$ are pairwise distinct elements of
$\{x_1,\ldots,x_n,y_1,\ldots,y_n\}$ and $z_i\neq x_i$ for all
$i=1,\ldots,n$.
%
We associate to $X$ a directed graph $G_X$ on the vertex set
$\{1,\ldots,n\}$ as follows: for $1\le i,j\le n$, there is a directed edge
from $i$ to $j$ if and only if $x_i=z_j$.
}
\end{Definition}

For the directed graph $G_X$ associated with the matrix $X$, directed loops cannot occur, since a loop at the vertex $i$ would mean $x_i=z_i$, contrary to the assumption $z_i\neq x_i$. Hence $G_X$ has no directed cycles of length $1$.

The following remark describes the structure of the connected components of the graph $G_X$.

\begin{Remark}\label{rem:def}
{\em The graph $G_X$ has a particularly simple structure. By construction, every
vertex of $G_X$ has indegree at most one. Since the entries
$z_1,\ldots,z_n$ are pairwise distinct, every vertex has outdegree at most
one. Hence every connected component of $G_X$ is either a directed path or
a directed cycle.}
\end{Remark}

The following example illustrates the construction of the graph $G_X$.

\begin{Example}\label{example}
{\em 
(1)
Consider the matrix
\[
X=
\begin{bmatrix}
x_1 & x_2 & x_3 & x_4 & x_5\\
x_2 & x_1 & y_1 & x_3 & y_5
\end{bmatrix}.
\]
Then $G_X$ has vertex set $\{1,2,3,4,5\}$ and directed edges
\[
1\to2,\qquad
2\to1,\qquad
3\to4,
\]
since $x_1=z_2$, $x_2=z_1$, and $x_3=z_4$. Thus $G_X$ has three connected
components: a directed cycle on the vertices $\{1,2\}$, a directed path $3\to4$, and the isolated vertex $5$.  If $I$ denotes the ideal generated by the $2$-minors of $X$, then
\[
x_1^2-x_2^2=(x_1-x_2)(x_1+x_2)\in I,
\]
while neither factor belongs to $I$. Hence $I$ is not prime. 

(2) Consider the matrix
\[
Y=
\begin{bmatrix}
x_1 & x_2 & x_3 & x_4 & x_5 & x_6\\
z_1 & z_2 & x_1 & z_4 & x_2 & x_3
\end{bmatrix},
\]
where $z_1,z_2,z_4$ are distinct variables in
$\{y_1,\ldots,y_6\}$. Then $G_Y$ has three connected components,
namely the directed paths
\[
1\to3\to6,\qquad
2\to5,
\]
and the isolated vertex $4$. Relabeling the variables by setting
\[
z_1=y_1,\qquad
z_2=y_3,\qquad
z_4=y_6,
\]
and then reordering the columns according to the connected components,
we obtain the matrix
\[
\begin{bmatrix}
x_1 & x_3 & x_6 & x_2 & x_5 & x_4\\
y_1 & x_1 & x_3 & y_3 & x_2 & y_6
\end{bmatrix}.
\]
Hence, after relabeling the variables and permuting the columns of $Y$,
the resulting matrix is precisely $X_6(\delta)$, where
$\delta=\{2,3,5\}$. Consequently, if $J$ denotes the ideal generated by
the $2$-minors of $Y$, then $J$ is identified with the ideal
$I_6(\delta)$.
}
\end{Example}

The next lemma isolates the obstruction to primeness arising from a
directed cycle and will be used in the proof of Theorem~\ref{prime}.

\begin{Lemma}\label{lemma:cycle}
Let $R=K[u_1,\ldots,u_r]$ be a polynomial ring and
\[
U=
\begin{bmatrix}
u_1 & u_2 & \cdots & u_r\\
u_2 & u_3 & \cdots & u_1
\end{bmatrix}.
\]
Then the ideal generated by the $2$-minors of $U$ is not prime.
\end{Lemma}


\begin{proof}
Let $J$ be the ideal generated by the $2$-minors of $U$. Throughout the proof, indices are taken modulo $r$, in particular, $u_{r+1}=u_1$. We first claim that
\[
u_1^r-u_2^r\in J.
\]
Indeed, for $k=1,\ldots,r$, the minor of the first and $k$-th columns gives
\begin{equation}\label{eq1}
u_1u_{k+1}-u_2u_k\in J.
\end{equation}
We prove by induction on $k$ that
\begin{equation}\label{eq:ind}
u_1^{k-1}u_{k+1}-u_2^k\in J
\end{equation}
for all $k=1,\ldots,r$.

For $k=1$, the assertion is trivial. Assume that $k\geq 2$ and that the assertion holds for $k-1$. Multiplying \eqref{eq1} by $u_1^{k-2}$ gives
\begin{equation}\label{eq2}
u_1^{k-1}u_{k+1}-u_1^{k-2}u_2u_k\in J.
\end{equation}
By the induction hypothesis, $u_1^{k-2}u_k-u_2^{k-1}\in J.$ Multiplying this by $u_2$, we obtain
\begin{equation}\label{eq3}
u_1^{k-2}u_2u_k-u_2^k\in J.
\end{equation}
Adding \eqref{eq2} and \eqref{eq3}, we get
\[
u_1^{k-1}u_{k+1}-u_2^k\in J.
\]
This proves \eqref{eq:ind}. Taking $k=r$ in \eqref{eq:ind}, and using $u_{r+1}=u_1$, we obtain
\[
u_1^r-u_2^r\in J.
\]
This proves the claim.

Now write
\[
u_1^r-u_2^r
=
(u_1-u_2)
\left(
u_1^{r-1}+u_1^{r-2}u_2+\cdots+u_2^{r-1}
\right).
\]
We next prove that
\[
u_1^{r-1}+u_1^{r-2}u_2+\cdots+u_2^{r-1}
\equiv
u_r^{r-2}(u_1+u_2+\cdots+u_r)
\pmod J.
\]
To prove this congruence, note that the minor of the $k$-th and $r$-th
columns yields $u_1u_k-u_ru_{k+1}\in J$ for $k=1,\ldots,r$. Hence
\begin{equation}\label{eq:colr}
u_1u_k
\equiv
u_ru_{k+1}
\pmod J
\qquad \text{for } k=1,\ldots,r.
\end{equation}
Similarly, the minor of the first and $k$-th columns yields
$u_1u_{k+1}-u_2u_k\in J$ for $k=1,\ldots,r$. Hence
\begin{equation}\label{eq:col1}
u_2u_k
\equiv
u_1u_{k+1}
\pmod J
\qquad \text{for } k=1,\ldots,r.
\end{equation}
We first claim that, for $i=0,\ldots,r-1$,
\begin{equation}\label{eq:claim}
u_1^{r-1-i}u_2^i
\equiv
u_1^{r-2}u_{i+1}
\pmod J.
\end{equation}
% Using (\ref{eq:colr}), we obtain
% \[
% u_1^{r-2} u_{i+1}
% =
% u_1^{r-3}(u_1 u_{i+1})
% =
% u_1^{r-3} u_r u_{i+2}.
% \]
% Applying (\ref{eq:colr}) again to the factor $u_1 u_{i+2}$, we get
% \[
% u_1^{r-3} u_r u_{i+2}
% =
% u_1^{r-4} u_r (u_1 u_{i+2})
% =
% u_1^{r-4} u_r^2 u_{i+3}
% \]
For $i=0$, the claim is immediate. Let $i\geq 1$. Repeated use of
\eqref{eq:col1} yields
\[
u_2^i
\equiv
u_1^{i-1}u_{i+1}
\pmod J.
\]
Multiplying by $u_1^{r-1-i}$, we obtain
\[
u_1^{r-1-i}u_2^i
\equiv
u_1^{r-2}u_{i+1}
\pmod J.
\]
Next, repeated application of \eqref{eq:colr} gives
\[
u_1^{r-2}u_{i+1}
\equiv
u_1^{r-2-t}u_r^t u_{i+1+t}
\pmod J
\]
for every $t=0,\ldots,r-2$. Taking $t=r-2$, we obtain
\[
u_1^{r-2}u_{i+1}
\equiv
u_r^{r-2}u_{i+r-1}
\pmod J.
\]
Since the indices are taken modulo $r$, we have
$u_{i+r-1}=u_{i-1}$. Therefore
\[
u_1^{r-2}u_{i+1}
\equiv
u_r^{r-2}u_{i-1}
\pmod J.
\]
Combining this with \eqref{eq:claim}, we obtain
\[
u_1^{r-1-i}u_2^i
\equiv
u_r^{r-2}u_{i-1}
\pmod J
\qquad \text{for } i=0,\ldots,r-1.
\]
As $i$ runs from $0$ to $r-1$, the indices $i-1$ run through all
residue classes modulo $r$. Hence the right-hand side runs through the
monomials
\[
u_r^{r-2}u_1,\,
u_r^{r-2}u_2,\,
\ldots,\,
u_r^{r-2}u_r.
\]
Summing over $i=0,\ldots,r-1$, we obtain
\[
u_1^{r-1}+u_1^{r-2}u_2+\cdots+u_2^{r-1}
\equiv
u_r^{r-2}(u_1+u_2+\cdots+u_r)
\pmod J.
\]
Since $u_1^r-u_2^r\in J$, it follows that
\[
(u_1-u_2)u_r^{r-2}(u_1+\cdots+u_r)\in J.
\]
However, $J$ is a homogeneous ideal generated by quadratic forms. Hence
$J$ contains no nonzero linear forms. In particular,
\[
u_1-u_2\notin J, \quad u_r\notin J
\qquad\text{and}\qquad
u_1+\cdots+u_r\notin J.
\]
If $J$ were prime, it would follow that one of
$u_1-u_2$, $u_r$, or $u_1+\cdots+u_r$ belongs to $J$, a contradiction.
Hence $J$ is not prime.
\hspace{15.1cm}
\end{proof}


\begin{Theorem}
\label{prime}
Let
\[
X=
\begin{bmatrix}
x_1 & x_2 & \cdots & x_n\\
z_1 & z_2 & \cdots & z_n
\end{bmatrix},
\]
where $z_1,\ldots,z_n$ are pairwise distinct elements of $\{x_1,\ldots,x_n,y_1,\ldots,y_n\}$ and $z_i\neq x_i$ for all $i=1,\ldots,n$. Then the ideal of $2$-minors of $X$ is prime if and only if $G_X$ contains no directed cycle.
\end{Theorem}

\begin{proof}
Suppose that $G_X$ contains a directed cycle $i_1\to i_2\to\cdots\to i_r\to i_1.$ Then, by definition,
\[
x_{i_1}=z_{i_2},\quad
x_{i_2}=z_{i_3},\quad
\ldots,\quad
x_{i_r}=z_{i_1}.
\]
Hence, after a cyclic permutation of its columns, the corresponding
submatrix of $X$ is
\[
\begin{bmatrix}
x_{i_1} & x_{i_2} & \cdots & x_{i_r}\\
x_{i_2} & x_{i_3} & \cdots & x_{i_1}
\end{bmatrix}.
\]
More precisely, the proof of Lemma~2.4 shows that
\[
(x_{i_1}-x_{i_2})x_{i_r}^{\,r-2}
(x_{i_1}+x_{i_2}+\cdots+x_{i_r})\in I.
\]
Since $I$ is a homogeneous ideal generated by quadratic forms, it
contains no nonzero linear forms. In particular,
\[
x_{i_1}-x_{i_2}\notin I,\qquad
x_{i_r}\notin I,\qquad\text{and}\qquad
x_{i_1}+\cdots+x_{i_r}\notin I.
\]
If $I$ were prime, it would follow that one of $x_{i_1}-x_{i_2}$, $x_{i_r}$ or $x_{i_1}+\cdots+x_{i_r}$ belongs to $I$, a contradiction. Hence $I$ is not prime.

Now, assume that $G_X$ contains no directed cycle. Then, by
Remark~\ref{rem:def}, every connected component of $G_X$ is a directed
path. Let $i_1\to i_2\to\cdots\to i_r$ be such a path. If $r\ge2$, then
\[
x_{i_1}=z_{i_2},\quad
x_{i_2}=z_{i_3},\quad
\ldots,\quad
x_{i_{r-1}}=z_{i_r}.
\]
Since $i_1$ has indegree zero, no entry among
$x_1,\ldots,x_n$ is equal to $z_{i_1}$. Hence
$z_{i_1}\in\{y_1,\ldots,y_n\}$. Therefore, after ordering the columns as
$i_1,i_2,\ldots,i_r$, the corresponding submatrix of $X$
\[
\begin{bmatrix}
x_{i_1} & x_{i_2} & \cdots & x_{i_r}\\
z_{i_1} & x_{i_1} & \cdots & x_{i_{r-1}}
\end{bmatrix},
\]
can be transformed into $X_r(\{2,\ldots,r\})$ by relabeling the variables.

Repeating this construction for each connected component, we obtain a
decomposition of $X$ into blocks of the form
$X_r(\{2,\ldots,r\})$ and single columns whose second-row entries are
$y$-variables. Equivalently, after a suitable permutation of the columns
and relabeling of the variables, the matrix $X$ is of the form
$X_n(\delta)$ for some $\delta\subseteq\{2,\ldots,n\}$. Therefore, the
ideal generated by the $2$-minors of $X$ is prime by (i) of 
Theorem~\ref{prime+CM+linear}.
\hspace{10.1cm}
\end{proof}

We can now state the main classification theorem.

\begin{Theorem}
\label{mainPCML}
Let $S=K[x_1, \ldots, x_n,y_1, \ldots,y_n]$ denote the polynomial ring in $2n$ variables over a field $K$ and $X$ a $2 \times n$ matrix
\[
X=
\begin{bmatrix}
x_1 & x_2 & \cdots & x_n\\
z_1 & z_2 & \cdots & z_n
\end{bmatrix},
\]
where $z_1,\ldots,z_n$ are pairwise distinct elements of $\{x_1,\ldots,x_n,y_1,\ldots,y_n\}$ and where $z_i\neq x_i$ for all $i=1,\ldots,n$.  Then the following conditions are equivalent:
\begin{enumerate}
\item[(i)] $X$ is a PCML matrix;
\item[(ii)] $I_2(X)$ is prime;
\item[(iii)] the directed graph $G_X$ contains no directed cycle.
\end{enumerate}
\end{Theorem}

\begin{proof}
The implication $(i)\Rightarrow(ii)$ follows immediately from the
definition of a PCML matrix, while $(ii)\Leftrightarrow(iii)$ is
Theorem~\ref{prime}.

It remains to prove $(iii)\Rightarrow(i)$. Suppose that $G_X$ contains
no directed cycle. By Remark~\ref{rem:def}, every connected component of $G_X$
is a directed path. As shown in the proof of Theorem~\ref{prime}, after a
suitable permutation of the columns and relabeling of the variables,
the matrix $X$ is of the form $X_n(\delta)$ for some $\delta\subseteq\{2,\ldots,n\}$. By Theorem~\ref{prime+CM+linear}, the ideal
$I_n(\delta)$ is prime, the quotient $S_\delta/I_n(\delta)$ is
Cohen--Macaulay, and $I_n(\delta)$ has a linear resolution.
These properties are preserved under relabeling of variables and
permutation of columns, and adjoining unused variables does not affect
them. Hence $X$ is a PCML matrix.
\end{proof}



\section*{Acknowledgments} The second author is supported by Scientific and Technological Research Council of Turkey T\"UB\.{I}TAK under the Grant No: 124F113, and is thankful to T\"UB\.{I}TAK for its support.




\begin{thebibliography}{99}


\bibitem{BrunsVetter}
W.~Bruns and U.~Vetter,
``Determinantal Rings,''
Lecture Notes in Mathematics, Vol.~1327,
Springer-Verlag, Berlin, 1988.



\bibitem{ConcaRNC}
A.~Conca,
Hilbert function and resolution of the powers of the ideal of the
rational normal curve,
\emph{J. Pure Appl. Algebra} \textbf{152} (2000), 65--74.




\bibitem{Eis}
D.~Eisenbud, Introduction to algebras with straightening laws, {\em in} ``Ring Theory and Algebra III'' (B.~R.~McDonald, Ed.), Lect. Notes in Pure and Appl. Math., Volume 55, Marcel Dekker, New York, 1980, pp. 243--268.



\bibitem{EisenbudSturmfels}
D.~Eisenbud and B.~Sturmfels,
Binomial ideals,
\emph{Duke Math. J.} \textbf{84} (1996), 1--45.



\bibitem{HHgtm260}
J.~Herzog and T.~Hibi, ``Monomial Ideals,'' GTM 260, Springer, 2011.



\bibitem{Hibi}
T.~Hibi, Distributive lattices, affine semigroup rings and algebras with straightening laws, {\it in} ``Commutative Algebra and Combinatorics'' (M. Nagata and H. Matsumura, Eds.), Advanced Studies in Pure Math., Volume 11, North--Holland, Amsterdam, 1987, pp. 93--109.


\bibitem{HQS}
T.~Hibi, A.~A.~Qureshi and S.~Saeedi Madani,
Prime and Cohen--Macaulay binomial ideals with linear resolution,
arXiv:2607.22127, 2026.


\end{thebibliography}
\bigskip

\begin{footnotesize}
{\bf Declaration of competing interest.}  The authors declare that they have no known competing financial interests or personal relationships that could have appeared to influence the work reported in this paper.

\medskip

{\bf Data availability.} No data was used for the research described in the article. 

\medskip

\end{footnotesize}



\end{document}
